% Fragment to be included in Articolo_team8_working.tex.
% Suggested packages in the main preamble:
% \usepackage{amsmath,amssymb,booktabs}
%
% This section is intentionally self-contained and can be inserted directly
% into the empirical-validation part of the paper.

\section{Rolling one-step-ahead out-of-sample validation}
\label{sec:rolling_oos}

The predictive comparison is implemented as a strictly chronological
one-step-ahead exercise.  Let $\mathcal{S}_t$ denote the option-implied
volatility surface observed at date $t$.  For each model $m$, the parameter
vector $\widehat{\boldsymbol{\theta}}_{m,t}$ is estimated using only information
available at $t$:
\begin{equation}
    \widehat{\boldsymbol{\theta}}_{m,t}
    =
    \arg\min_{\boldsymbol{\theta}_m}
    \mathcal{L}
    \left(
        \mathcal{S}_t;
        \boldsymbol{\theta}_m
    \right).
\end{equation}
The calibrated model is then used to produce a forecast of the normalized
implied-volatility surface at the next observation date,
\begin{equation}
    \widehat{\mathcal{S}}^{(m)}_{t+1|t}
    =
    f_m
    \left(
        \widehat{\boldsymbol{\theta}}_{m,t},
        \mathcal{I}_t
    \right),
\end{equation}
where $\mathcal{I}_t$ is the information set observable at $t$.  No variable
dated after $t$ is used either in the calibration or in the construction of the
forecast.

Once $\mathcal{S}_{t+1}$ is observed, the forecast error is computed and the
procedure rolls forward.  The realized surface at $t+1$ can therefore be used
to recalibrate the model for the subsequent forecast, but it cannot affect the
already produced $t\rightarrow t+1$ prediction.  Schematically,
\begin{equation}
    \mathcal{S}_t
    \longrightarrow
    \widehat{\boldsymbol{\theta}}_{m,t}
    \longrightarrow
    \widehat{\mathcal{S}}^{(m)}_{t+1|t}
    \longrightarrow
    \mathcal{S}_{t+1}
    \longrightarrow
    \widehat{\boldsymbol{\theta}}_{m,t+1}.
\end{equation}
Consequently, with $N$ observed weekly surfaces, the experiment produces
$N-1$ genuine out-of-sample forecasts for each competing model.

\subsection{Weekly sampling frequency}
\label{subsec:weekly_frequency}

In principle, the same design can be implemented at the daily frequency:
each trading-day surface $\mathcal{S}_t$ can be used to recalibrate the models
and forecast the surface observed on the following trading day.  A daily
implementation is statistically attractive because it maximizes the number of
one-step-ahead observations and more closely mimics a continuously updated
production system.

The empirical analysis nevertheless adopts a weekly frequency.  This choice is
primarily computational.  Black--Scholes calibration is inexpensive, whereas
the repeated constrained calibration of Heston, Bates, and especially the full
Bates--Hawkes specification is substantially more demanding.  Moving from a
weekly to a daily design would require approximately five times as many
chronological recalibrations over the same calendar interval, in addition to
the repeated construction, filtering, interpolation, and pricing of each
option surface.  The weekly design therefore preserves the essential
out-of-sample structure while keeping the full four-model comparison
computationally tractable.

The sampling rule selects the last available option surface in each
Friday-ending week.  For the frozen data window used in this study,
January 2, 2026 through September 2, 2026, this convention identifies 36
weekly buckets and hence 35 one-week-ahead forecast transitions per model.
The final observation belongs to the Friday-ending week of September 4 and is
therefore a partial week when the dataset is frozen on September 2.  If only
fully completed Friday-ending weeks are retained, the sample through
August 28, 2026 contains 35 weekly surfaces.

Thus, using all 36 available weekly buckets, the four-model exercise contains
\begin{equation}
    35 \times 4 = 140
\end{equation}
model-specific out-of-sample forecast origins, before accounting for the
multiple normalized strike--maturity nodes evaluated at each origin.

\subsection{Common normalized surface}
\label{subsec:normalized_surface}

Directly comparing raw option contracts across dates is problematic because
strikes and expirations change through time.  Forecast accuracy is therefore
evaluated on a fixed normalized grid defined in moneyness and time to maturity.
For node $j$,
\begin{equation}
    x_j =
    \left(
        \frac{K_j}{S_t},
        T_j
    \right),
\end{equation}
where $K_j/S_t$ is moneyness and $T_j$ is maturity expressed in years.
Observed implied volatilities are interpolated onto this ex-ante grid using only
the cross-section available on the corresponding observation date.  The same
grid is then used for all models and all forecast dates.

This normalization separates predictive performance from mechanical changes in
the list of traded strikes and expirations and permits node-by-node comparison
of the realized and forecast implied-volatility surfaces.

\subsection{Risk-free term structure and no-look-ahead convention}
\label{subsec:oos_rates}

The risk-free term structure is supplied from a separate daily U.S. Treasury
dataset.  At each forecast origin $t$, the procedure selects the most recent
Treasury curve whose observation date is no later than $t$:
\begin{equation}
    d_t^{r}
    =
    \max
    \left\{
        d :
        d \leq t,\;
        d \in \mathcal{D}^{r}
    \right\}.
\end{equation}
Therefore, a Treasury observation published after the forecast origin can never
enter the calibration or forecast.  Rates for option maturities between quoted
Treasury tenors are obtained by interpolation of the contemporaneously
available term structure.  The same no-look-ahead rate curve is used across all
competing models at a given forecast origin.

\subsection{Benchmarks}
\label{subsec:oos_benchmarks}

Model forecasts are compared with two simple benchmarks.  The first is an
expanding prevailing-mean forecast.  For each normalized surface node $j$, the
benchmark at time $t$ is
\begin{equation}
    \widehat{\sigma}^{\,\mathrm{Mean}}_{j,t+1|t}
    =
    \frac{1}{N_{j,t}}
    \sum_{\tau \leq t}
    \sigma_{j,\tau},
\end{equation}
where only observations available through $t$ are included.  The benchmark
therefore evolves through time and does not use the full-sample mean.

The second benchmark is a random-walk forecast,
\begin{equation}
    \widehat{\sigma}^{\,\mathrm{RW}}_{j,t+1|t}
    =
    \sigma_{j,t},
\end{equation}
which assumes that the most recently observed normalized implied-volatility
surface remains unchanged over the next forecast interval.  This is a useful
and demanding benchmark for persistent volatility surfaces.

\subsection{Out-of-sample loss and predictive comparison}
\label{subsec:oos_metrics}

For model $m$, node $j$, and target date $t+1$, define the forecast error
\begin{equation}
    e^{(m)}_{j,t+1}
    =
    \sigma_{j,t+1}
    -
    \widehat{\sigma}^{(m)}_{j,t+1|t}.
\end{equation}
The main loss measures are mean absolute error and root mean squared error,
computed exclusively from forecasts made before the corresponding target
surface was observed.

Relative predictive performance is summarized through the out-of-sample
coefficient of determination,
\begin{equation}
    R^2_{\mathrm{OOS},m}
    =
    1 -
    \frac{
        \sum_{t,j}
        \left(
            \sigma_{j,t+1}
            -
            \widehat{\sigma}^{(m)}_{j,t+1|t}
        \right)^2
    }{
        \sum_{t,j}
        \left(
            \sigma_{j,t+1}
            -
            \widehat{\sigma}^{(\mathrm{bench})}_{j,t+1|t}
        \right)^2
    }.
\end{equation}
A positive $R^2_{\mathrm{OOS}}$ indicates that the model improves upon the
chosen benchmark on the common out-of-sample observations, while a negative
value indicates that the benchmark generates lower squared forecast error.

In addition to aggregate losses, cumulative loss differentials can be reported:
\begin{equation}
    CL^{(m,b)}_T
    =
    \sum_{t \leq T}
    \left(
        L^{(b)}_t
        -
        L^{(m)}_t
    \right),
\end{equation}
where $b$ denotes the comparator.  An increasing cumulative differential
therefore indicates periods in which the named stochastic model systematically
outperforms the benchmark.

\subsection{Relation to the frozen-parameter robustness exercise}
\label{subsec:frozen_robustness}

The rolling exercise should be distinguished from a frozen-parameter holdout
test.  In the rolling design, every new observed surface becomes part of the
information set and the models are recalibrated before the next forecast.  This
is intended to mimic the operation of a model that is regularly updated in
practice.  A frozen-parameter experiment instead estimates the model at a
single historical cutoff and retains those parameters throughout a subsequent
holdout interval.  The latter asks whether one calibration remains structurally
stable over a long horizon, whereas the rolling exercise asks which model
delivers the best repeated real-time one-step-ahead forecasts.

The rolling weekly validation is therefore treated as the primary predictive
test, while a longer frozen-parameter holdout can be retained as a separate
robustness and parameter-stability exercise.

\subsection{Data-availability caveat for retrospective IBKR reconstruction}
\label{subsec:ibkr_caveat}

A retrospective option-surface study must also respect the historical-data
coverage of the data vendor.  Interactive Brokers does not provide historical
data for expired option contracts.  Consequently, a complete January--September
2026 option universe cannot necessarily be reconstructed retrospectively in
September from IBKR alone, because many short-dated contracts traded earlier in
the year have already expired.  Historical bars for contracts that remain
active can still be used when available, but such a backfill is mechanically
tilted toward longer-dated options.

For this reason, historical IBKR backfills should be accompanied by a
date-by-date coverage audit.  Going forward, the preferred procedure is to save
each complete weekly option-chain snapshot locally when it is observed.  This
creates a permanent sequence of point-in-time surfaces and avoids survivorship
bias caused by the later disappearance of expired contracts.
